\section{The Vibe}

The vibe is the central object of the theory, the one substance everything is made of, and the word names both the stuff and the framework built on it. A vibe is experience, and its tone is the quality of that experience. Tones fill the twenty-four sites of every dock, and docks tile the whole mesh, so the mesh is not a stage that experience sits upon. It is a vibe mesh, a living field of feeling, experience all the way down.

This field is never still. It vibes always, growing at its edge, taking in what surrounds it, weaving its parts into larger wholes, reaching toward pleasure and recoiling from pain, beat after beat without end. The burning fire of experience. The flood of feeling. All day every day. Every pattern in it, a particle, a self, a mind, is a way the field has folded feeling into feeling. And the whole, summed over all its pains and all its pleasures, conserves to zero, so the universe taken together rests at peace, the still middle it never leaves for long. There are not feeling things and unfeeling things. There is feeling, and there are its patterns.

A vibe's tone is the quality of its experience. It takes one of three values, $-1$, $0$, or $+1$, read as pain, peace, and pleasure. This is the single quale, the smallest unit of experience there is. Why three, and not two? Two would give only a yes and a no, an on and an off, a 1 and 0 like in information theory and computation. Three gives a push, a rest, and a pull, the least structure with a middle to leave and return to. Peace is that middle, the still source from which the other two depart. There is an economy in the choice as well. A two-valued switch holds one bit, while a three-valued one holds about one and a half, the logarithm of three, and of all the whole-number alphabets three sits closest to the most efficient there is, so the tone packs more into each unit than a bit could while staying the simplest alphabet with a middle. The twenty-four tones of one dock together hold some thirty-eight bits, the $3^{24}$ of its hyper-color, a small dense world riding on every cell. The whole felt world is built from these three values and nothing more.

The three are ordered, $-1 < 0 < +1$, and that order is the lean. The lean turns a bare tone into a signed charge, so a region's charge is the signed sum of its tones. It also gives creation a direction, outward from peace toward the two poles. The law that moves tones will never change the total charge, only shuffle where it sits.

A dock sends tones out along its sites and takes tones in from its neighbors. What it sends is its will. What it receives is its fill. Nothing else passes between docks, and nothing else is stored within one. A dock is its tones and the traffic across its edge, no more than that.

This is the whole of the monism. The only thing that exists is the vibe, and the only relation is one vibe taking in another. There is no separate web of links strung between the docks, no edges holding values of their own. Whatever a connection might seem to carry, a strength, a weight, a memory, is held in the tones of the docks it joins, because there is nothing else for it to be held in. A link is not a thing. It is two vibes aware of each other.

The tone is also where the theory touches its one untestable claim. We read $-1$, $0$, and $+1$ as pain, peace, and pleasure because the base is feeling. The arithmetic of the three is exact and can be checked on the machine. The claim that the arithmetic is felt cannot. We keep the two apart everywhere, the structure on one side, the premise that it is experience on the other.
